% TEMPLATE-MILC-v3.tex - plantilla maestra UNICA de las guias MILC v3.
%
% La usan las tres familias de guias, cada una con su driver:
%   · Grados 6-11  -> scripts/build-guias-g11.py
%   · Bebras       -> scripts/build-guias-bebras.py
%   · Semillero    -> scripts/build-guias-semillero.py
%
% Antes habia tres copias casi identicas (~400 lineas iguales, 6 distintas):
% cada mejora habia que aplicarla tres veces, y en la practica no se hacia
% (el motor de recursos vivia solo en dos de ellas). Lo que varia entre
% familias esta parametrizado en 6 placeholders que cada driver llena
% (se nombran aqui SIN su sintaxis real: el builder reemplaza por texto plano,
% tambien dentro de los comentarios, y un valor multilinea rompe el preambulo):
%
%   HEADER_IZQ        encabezado izquierdo de pagina
%   HEADER_DER        encabezado derecho de pagina
%   PORTADA_ETIQUETA  rotulo sobre el numero grande (GRADO/MOMENTO/MODULO)
%   PORTADA_SERIE     linea de serie de la portada
%   PORTADA_ID        identificador de la guia en la portada
%   PORTADA_CIRCULO   el circulito de grado (°), o vacio si no aplica
%
% El resto de claves las documenta content/guias/_SCHEMA.md. El script valida
% que no queden placeholders sin reemplazar antes de escribir el archivo final.

\documentclass[11pt,letterpaper]{article}
\usepackage[margin=2.0cm]{geometry}
\usepackage[table]{xcolor}
\usepackage{fontspec}
\usepackage[spanish]{babel}
\usepackage{tikz}
% Librerías TikZ para los diagramas de las guías (bloque `recursos:`):
%  · babel  → babel en español vuelve activos caracteres como «>», y TikZ los usa
%             en sus opciones (>=stealth, ->). Sin esta librería, cualquier
%             diagrama con flechas revienta con «\language@active@arg> has an
%             extra }». Debe ir ANTES de que se dibuje nada.
%  · positioning        → sintaxis «below left=2pt and 3pt».
%  · decorations.pathreplacing → llaves ({brace}) para señalar rangos.
\usetikzlibrary{babel,positioning,decorations.pathreplacing}
\usepackage{graphicx}
% Circuitos: el trabajo de electrónica se hace hoy con micro:bit y sus sensores
% integrados (MakeCode, sin cableado), así que no se necesitan esquemáticos.
% Si algún día entran montajes con protoboard, basta descomentar esta línea:
% los diagramas .tex/.tikz declarados en `recursos:` ya se incrustan solos.
% \usepackage{circuitikz}
\usepackage[most]{tcolorbox}
\usepackage{tabularx}
\usepackage{array}
\usepackage{fancyhdr}
\usepackage{titlesec}
\usepackage[document]{ragged2e}  % bandera a la derecha en todo el cuerpo
\usepackage{enumitem}
\usepackage{microtype}

% Helvetica Neue no trae la flecha U+2192, asi que cada «→» escrito literalmente
% en el YAML se compilaba a NADA, en silencio (462 flechas en 61 guias: rutas del
% tipo «paleta → tipografia → lockup» salian rotas). Mapeamos el caracter a su
% version matematica, que si tiene glifo. Asi el autor puede seguir escribiendo
% «→» comodamente en el YAML.
\usepackage{newunicodechar}
\newunicodechar{→}{\ensuremath{\rightarrow}}
% Mismo problema con los simbolos de marca y vinetas: el «✓» de «una palomita ✓»
% salia como cuadrito vacio en el PDF de todas las guias que lo usaban.
\usepackage{pifont}
\newunicodechar{✓}{\ding{51}}
\newunicodechar{✗}{\ding{55}}
\newunicodechar{✔}{\ding{52}}
\newunicodechar{•}{\textbullet}
\newunicodechar{≥}{\ensuremath{\geq}}
\newunicodechar{≤}{\ensuremath{\leq}}

\setmainfont{Helvetica Neue}
\setsansfont{Helvetica Neue}
\setlength{\parindent}{0pt}
\setlength{\parskip}{6pt}
% Sin particion de palabras: en 6.º-7.º una palabra cortada por guion al final
% de linea («Comuni-cacion») frena la lectura mucho mas que una linea corta.
\hyphenpenalty=10000
\exhyphenpenalty=10000
\pretolerance=2000
\tolerance=3000
\setlength{\emergencystretch}{3em}
\linespread{1.10}
\renewcommand{\arraystretch}{1.25}
\setlist{leftmargin=5mm,itemsep=1mm,topsep=1mm}
\newcolumntype{Y}{>{\RaggedRight\arraybackslash}X}

\definecolor{milcMagenta}{HTML}{E6007E}
\definecolor{milcVino}{HTML}{5A0038}
\definecolor{milcTurquesa}{HTML}{0093A5}
\definecolor{milcVerde}{HTML}{58A923}
\definecolor{milcMostaza}{HTML}{E5B400}
\definecolor{milcOcre}{HTML}{B86600}
\definecolor{milcNegro}{HTML}{241816}
\definecolor{milcGris}{HTML}{F7F7F4}
\definecolor{milcLinea}{HTML}{E8E2DD}
\definecolor{milcGradePrimary}{HTML}{0066FF}
\definecolor{milcGradeDark}{HTML}{003D99}
\definecolor{milcGradeSoft}{HTML}{D6E8FF}
\definecolor{milcGradeAccent}{HTML}{CFE7FF}
\definecolor{milcGradeText}{HTML}{FFFFFF}
\color{milcNegro}

\pagestyle{fancy}
\fancyhf{}
\lhead{\small Tecnología e Informática · Grado 6}
\rhead{\small Guía 12 · MILC v3}
\cfoot{\small\thepage}
\renewcommand{\headrulewidth}{0pt}

\newtcolorbox{titlebox}[1]{
  enhanced,colback=#1,colframe=#1,arc=7mm,boxrule=0pt,
  left=7mm,right=7mm,top=3mm,bottom=3mm,
  before skip=8pt,after skip=8pt,
  fontupper=\sffamily\bfseries\Large\color{white}
}

\newtcolorbox{softbox}[3]{
  enhanced,colback=#2,colframe=#1,borderline west={4pt}{0pt}{#1},
  arc=2mm,boxrule=.4pt,left=4mm,right=4mm,top=3mm,bottom=3mm,
  before skip=6pt,after skip=8pt,title={#3},coltitle=white,
  colbacktitle=#1,fonttitle=\sffamily\bfseries,fontupper=\RaggedRight,
  attach boxed title to top left={xshift=3mm,yshift=-2mm},
  boxed title style={arc=2mm, boxrule=0pt}
}

\newtcolorbox{drawbox}[2]{
  enhanced,colback=white,colframe=#1,arc=4mm,boxrule=1pt,
  left=5mm,right=5mm,top=4mm,bottom=4mm,before skip=8pt,after skip=8pt,
  title={#2},coltitle=white,colbacktitle=#1,fonttitle=\sffamily\bfseries,
  fontupper=\RaggedRight,
  attach boxed title to top left={xshift=4mm,yshift=-2mm},
  boxed title style={arc=3mm, boxrule=0pt}
}

\newtcolorbox{infoband}[1]{
  enhanced,colback=#1!8,colframe=#1!50,arc=2mm,boxrule=.5pt,
  left=4mm,right=4mm,top=2mm,bottom=2mm,before skip=4pt,after skip=4pt,
  fontupper=\small\RaggedRight
}

% Puente narrativo entre secciones (contrato editorial Conéctate).
% Da continuidad: "Acabas de X. Ahora vas a Y porque Z."
% Visualmente: banda izquierda en color del grado, texto en itálica pequeña.
\newtcolorbox{puentebox}{
  enhanced,colback=milcGradeSoft,colframe=milcGradeDark,
  borderline west={3pt}{0pt}{milcGradeDark},
  arc=0mm,boxrule=0pt,left=5mm,right=4mm,top=2.5mm,bottom=2.5mm,
  before skip=4pt,after skip=8pt,
  fontupper=\itshape\small\color{milcNegro}\RaggedRight
}
% La flecha va en modo matematico a proposito: Helvetica Neue NO tiene el glifo
% U+2192, asi que \textrightarrow se compilaba a NADA («Missing character: There
% is no → in font Helvetica Neue») y el puente salia sin flecha. En modo
% matematico la flecha viene de la fuente matematica y si se dibuja.
\newcommand{\puente}[1]{%
  \begin{puentebox}%
    \textcolor{milcGradeDark}{$\rightarrow$}\hspace{2mm}#1%
  \end{puentebox}%
}

\newcommand{\linead}[1]{\noindent\color{milcLinea}\rule{#1}{0.5pt}\color{milcNegro}}
\newcommand{\respuesta}[1]{\par\vspace{2mm}\foreach \n in {1,...,#1}{\noindent\color{milcLinea}\rule{\linewidth}{0.5pt}\par\vspace{5mm}}\color{milcNegro}}
\newcommand{\checkbox}{\textcolor{milcGradeDark}{\fbox{\phantom{x}}}\hspace{5pt}}

\newcommand{\phasecard}[4]{
\begin{tcolorbox}[enhanced,colback=#3,colframe=#2,arc=3mm,boxrule=.5pt,height=3.05cm,valign=center,halign=center,left=1.5mm,right=1.5mm,top=2mm,bottom=2mm]
{\sffamily\bfseries\fontsize{22}{24}\selectfont\textcolor{#2}{#1}}\par
{\sffamily\bfseries\small #4}
\end{tcolorbox}
}

% ── Piezas del rediseno 2026 (legibilidad 6.º-11.º) ───────────────────
% Banda de estacion: reemplaza el titlebox macizo. Titulo a la izquierda,
% dato practico (tiempo/modalidad) a la derecha, en una sola linea.
\newcommand{\milcestacion}[2]{%
\begin{tcolorbox}[enhanced,colback=#1,colframe=#1,arc=2mm,boxrule=0pt,
  left=5mm,right=5mm,top=2.6mm,bottom=2.6mm,before skip=11pt,after skip=7pt]
{\sffamily\bfseries\fontsize{15}{18}\selectfont\color{white}#2}
\end{tcolorbox}}

% Tarjeta de paso numerada: sustituye las tablas «Pilar / Aplicacion», que
% eran formato de consulta docente, no de estudiante. El numero va en un
% circulo sobre el borde para que la secuencia se lea de un vistazo.
\newcommand{\milcpaso}[3]{%
\begin{tcolorbox}[enhanced,colback=white,colframe=milcLinea,arc=1.5mm,boxrule=.9pt,
  left=14mm,right=4mm,top=3mm,bottom=3mm,before skip=4pt,after skip=4pt,
  fontupper=\RaggedRight,
  overlay={\node[anchor=north west,circle,fill=#2,text=white,inner sep=0pt,
    minimum size=7.4mm,font=\sffamily\bfseries\small]
    at ([xshift=3.6mm,yshift=-3.2mm]frame.north west) {#1};}]
#3
\end{tcolorbox}}

% Casilla marcable de tamano util con lapiz (la anterior era \fbox{\phantom{x}},
% de alto variable segun el texto que la seguia).
\newcommand{\milcmarca}{\framebox[3.8mm]{\rule{0pt}{3.4mm}}}

% Fila marcable de lista suelta (checks de la estacion 1).
\newcommand{\milccheckrow}[1]{%
\par\vspace{1.8mm}\noindent
\makebox[6.4mm][l]{\raisebox{-.55mm}{\textcolor{milcNegro}{\milcmarca}}}%
\parbox[t]{\dimexpr\linewidth-6.4mm\relax}{\RaggedRight #1}\par}

% Celda marcable dentro de la rubrica (tabla de autoevaluacion). Misma
% sangria francesa, pero pensada para vivir dentro de una columna X.
\newcommand{\milcopcion}[1]{%
\makebox[5.2mm][l]{\raisebox{-.55mm}{\textcolor{milcNegro}{\milcmarca}}}%
\parbox[t]{\dimexpr\linewidth-5.2mm\relax}{\RaggedRight #1}}

% ── Recursos visuales de la guía ──────────────────────────────────────
% Declarados en el bloque `recursos:` del YAML y emitidos por el builder.
% \guiaFigura   → imagen rasterizada o PDF (foto, carta celeste, montaje).
% \guiaDiagrama → fuente TikZ/circuitikz (.tex) incrustada como vector:
%                 es la vía para circuitos y esquemáticos editables.
% #1 = ruta relativa al .tex · #2 = pie (puede ir vacío)
\newcommand{\guiaFigura}[2]{%
\begin{center}
\includegraphics[width=0.88\linewidth,height=0.38\textheight,keepaspectratio]{#1}\\[1.5mm]
{\footnotesize\itshape\textcolor{milcNegro!75}{#2}}
\end{center}
}

\newcommand{\guiaDiagrama}[2]{%
\begin{center}
\input{#1}\\[1.5mm]
{\footnotesize\itshape\textcolor{milcNegro!75}{#2}}
\end{center}
}

\begin{document}
\thispagestyle{empty}

% ── Portada ───────────────────────────────────────────────────────────
% Antes: un rectangulo de color a pagina completa. Se veia bien en pantalla,
% pero en la fotocopiadora del colegio se comia el toner y salia gris sucio.
% Ahora la banda ocupa el tercio superior y el resto va en blanco.
\begin{tikzpicture}[remember picture,overlay]
  \path[fill=milcGradePrimary]
    (current page.north west) rectangle ([yshift=-10.6cm]current page.north east);

  \node[anchor=north west,text=milcGradeText] at ([xshift=2.0cm,yshift=-1.55cm]current page.north west)
    {\sffamily\bfseries\fontsize{12}{14}\selectfont GRADO};
  \node[anchor=north west,text=milcGradeText,opacity=.85] at ([xshift=2.0cm,yshift=-2.35cm]current page.north west)
    {\sffamily\bfseries\fontsize{8.5}{10}\selectfont SERIE GUÍAS · TECNOLOGÍA E INFORMÁTICA};
  \node[anchor=north east,text=milcGradeText,opacity=.85,align=right] at ([xshift=-2.0cm,yshift=-1.55cm]current page.north east)
    {\sffamily\bfseries\fontsize{9}{12}\selectfont I.E. Sor María Juliana\\Cartago, Valle del Cauca};

  \node[anchor=west,text=milcGradeText] (gradonum) at ([xshift=1.85cm,yshift=-7.1cm]current page.north west)
    {\sffamily\bfseries\fontsize{112}{112}\selectfont 6};
  % El circulito de grado (°) solo aplica a las guias de grado («11°»); en
  % Bebras y Semillero el numero grande es un momento/modulo, no un grado.
  % Se posiciona relativo al nodo `gradonum`, no en coordenadas absolutas.
  \draw[milcGradeText,line width=1.6pt]
    ([xshift=2.0mm,yshift=-4.5mm]gradonum.north east) circle (.15cm);

  \node[anchor=west,text=milcGradeText,text width=11.4cm,align=left] at ([xshift=1.5cm,yshift=0.5cm]gradonum.east)
    {
     {\sffamily\bfseries\fontsize{9}{11}\selectfont\MakeUppercase{Período 2 · Hardware y software}}\\[3mm]
     {\sffamily\bfseries\fontsize{21}{25}\selectfont\setlength{\baselineskip}{27pt}¿Qué es un\\[4pt]computador? ---\\[4pt]las 4 funciones\\[4pt]de toda máquina\\[4pt]inteligente}\\[3.5mm]
     {\sffamily\bfseries\fontsize{15}{18}\selectfont Guía 12-6-TIC}
    };
\end{tikzpicture}

\vspace*{9.4cm}
\vfill

{\sffamily\bfseries\fontsize{8}{10}\selectfont\textcolor{milcGradeDark}{LO QUE TE LLEVAS HOY}}

\begin{tcolorbox}[enhanced,colback=white,colframe=milcNegro,arc=2mm,boxrule=1.6pt,
  left=5mm,right=5mm,top=4mm,bottom=4mm,before skip=3pt,after skip=12pt,
  fontupper=\RaggedRight\fontsize{11}{15.5}\selectfont]
Un \textbf{diagrama del computador con las 4 funciones etiquetadas} (entrada, proceso, almacenamiento, salida), hecho en una hoja del cuaderno. El diagrama incluye \textbf{flechas entre las funciones} mostrando el flujo de la información, \textbf{2 ejemplos concretos por cada función} (10 ejemplos en total contando una función adicional), y \textbf{un caso completo aplicado}: el recorrido de la información cuando juegas Roblox, desde tu dedo en la pantalla hasta el sonido que sale del altavoz.
\end{tcolorbox}

\begin{tcolorbox}[enhanced,colback=milcGris,colframe=milcGris,arc=2mm,boxrule=0pt,
  left=5mm,right=5mm,top=3.5mm,bottom=3.5mm,after skip=12pt]
{\sffamily\bfseries\fontsize{8}{10}\selectfont\textcolor{milcNegro!55}{ESTA GUÍA ES DE}}\\[3mm]
\begin{tabularx}{\linewidth}{X p{3.2cm} p{3.2cm}}
\linead{\linewidth} & \linead{3.0cm} & \linead{3.0cm}\\[-1mm]
{\fontsize{8.5}{10}\selectfont\textcolor{milcNegro!55}{Nombre completo}} &
{\fontsize{8.5}{10}\selectfont\textcolor{milcNegro!55}{Grupo}} &
{\fontsize{8.5}{10}\selectfont\textcolor{milcNegro!55}{Fecha}}\\
\end{tabularx}
\end{tcolorbox}

\vfill

\noindent\textcolor{milcLinea}{\rule{\linewidth}{1pt}}\\[2mm]
{\sffamily\bfseries\fontsize{9.5}{12}\selectfont Autor: Dr. Álvaro Cárdenas Orozco}\\[1mm]
{\sffamily\fontsize{8.5}{11}\selectfont\textcolor{milcNegro!55}{Metodología MILC · apertura ancestral · pensamiento computacional · triángulo Dussel–estoicismo–Floridi}}

\newpage

\section*{Apertura: ¿Qué es un computador? --- las 4 funciones básicas de toda máquina inteligente}

\begin{softbox}{milcGradeDark}{milcGradeSoft}{Saber ancestral}
\textbf{En cada barrio de Cartago vivía un relojero.} Hasta hace pocas décadas, en una vitrina pequeña de la calle 12 o la calle 14 trabajaba el relojero. Don Lucho, don Aurelio, doña Rita --- cada barrio tenía el suyo. \textbf{El relojero no era un mago: era un sabio de las máquinas pequeñas}. Sabía exactamente qué pieza hacía qué dentro del reloj. Conocía los engranajes que recibían la cuerda (\emph{entrada}), las ruedas que la transmitían (\emph{proceso}), el resorte que guardaba la energía (\emph{almacenamiento}), y las manecillas que mostraban la hora (\emph{salida}). Si una pieza fallaba, sabía cuál era. Si el reloj se atrasaba, sabía cuál engranaje ajustar. Las personas le tenían respeto porque \textbf{conocía la máquina por dentro}. Hoy el reloj de tu abuela ya casi no existe; hoy todos cargamos otra máquina por dentro: el computador y el celular. Aprender qué hay dentro te convierte en algo parecido al relojero: alguien que entiende la máquina, no alguien que solo la usa.
\end{softbox}

\begin{softbox}{milcTurquesa}{milcGris}{Saber contemporáneo}
Un \textbf{computador} es una máquina que hace siempre lo mismo: recibe datos, los procesa, los guarda y muestra resultados. Esas son las \textbf{4 funciones básicas}. Las hace tu computador del aula, las hace tu celular, las hace una caja registradora del supermercado, las hace un cajero electrónico. Cualquier máquina inteligente sigue ese mismo patrón. \textbf{Función 1 --- Entrada:} la información que entra (lo que escribes en el teclado, lo que dices al micrófono, lo que tocas en la pantalla). \textbf{Función 2 --- Proceso:} lo que el cerebro de la máquina (la CPU) hace con esa información. \textbf{Función 3 --- Almacenamiento:} dónde la guarda (memoria temporal --- RAM --- o memoria permanente --- disco duro). \textbf{Función 4 --- Salida:} cómo te muestra el resultado (pantalla, altavoz, impresora). Aprender estas 4 funciones es el primer paso para entender por dentro cualquier máquina.
\end{softbox}

\begin{softbox}{milcMostaza}{milcGris}{Pregunta-puente}
\textit{Cuando juegas Roblox en el celular y mueves tu personaje con el dedo en la pantalla, \textbf{¿qué partes del celular están trabajando en ese instante}? ¿Sabrías nombrar al menos 3?}
\end{softbox}

\begin{softbox}{milcVerde}{milcGris}{Saber y saber hacer}
Al terminar la clase: (1) podrás \textbf{identificar} las 4 funciones básicas de un computador; (2) sabrás \textbf{explicar} con tus palabras qué hace cada una; (3) podrás \textbf{aplicar} las 4 funciones a un caso real (un juego, un programa, una acción cotidiana); (4) habrás \textbf{creado} un diagrama propio del computador con las 4 funciones etiquetadas.
\end{softbox}



\small
\begin{tabularx}{\textwidth}{>{\bfseries}p{3.0cm}Y}
\rowcolor{milcGradePrimary!12}Autor & Dr. Álvaro Cárdenas Orozco\\
DBA & Identifica componentes y sistemas tecnológicos en su entorno (MEN, T\&I 6°)\\
Referentes & Saber ancestral del relojero · Sistemas como cadena de oficios · Diagnóstico paso a paso\\
Producto final & Un \textbf{diagrama del computador con las 4 funciones etiquetadas} (entrada, proceso, almacenamiento, salida), hecho en una hoja del cuaderno. El diagrama incluye \textbf{flechas entre las funciones} mostrando el flujo de la información, \textbf{2 ejemplos concretos por cada función} (10 ejemplos en total contando una función adicional), y \textbf{un caso completo aplicado}: el recorrido de la información cuando juegas Roblox, desde tu dedo en la pantalla hasta el sonido que sale del altavoz.\\
\end{tabularx}
\normalsize

\puente{Hoy haces 4 cosas: clasificas 10 acciones cotidianas según función, aprendes las 4 funciones con sus piezas, aplicas todo a un juego que conoces, y dibujas tu diagrama.}

\milcestacion{milcGradeDark}{Ruta MILC: así vamos a aprender}
\begin{tabularx}{\textwidth}{>{\centering\arraybackslash}X>{\centering\arraybackslash}X>{\centering\arraybackslash}X>{\centering\arraybackslash}X}
\phasecard{1}{milcVerde}{milcGris}{Escucha\\[-1mm]\small Observo, escucho y pregunto.} &
\phasecard{2}{milcTurquesa}{milcGris}{Entender\\[-1mm]\small Sistematizo y ordeno.} &
\phasecard{3}{milcMagenta}{milcGris}{Hacer\\[-1mm]\small Construyo y aplico.} &
\phasecard{4}{milcVino}{milcGris}{Cierre\\[-1mm]\small Evalúo y reflexiono.}
\end{tabularx}


\puente{Empezamos con un mini-cuestionario: 10 acciones cotidianas para clasificar en 4 funciones.}

\milcestacion{milcVerde}{Estación 1 de 3 · Escucha}

\begin{softbox}{milcVerde}{milcGris}{Escena para escuchar}
\textbf{Actividad 1 · IDENTIFICA --- 10 acciones para clasificar} (10 min · individual). En el cuaderno, dibuja una tabla de 2 columnas. Encabezado izquierda: \textbf{Acción}. Encabezado derecha: \textbf{¿Qué función es?} (entrada / proceso / almacenamiento / salida). Clasifica estas \textbf{10 acciones} que pasan cuando usas un computador o celular: (1) \emph{tocar la pantalla del celular para abrir TikTok}; (2) \emph{el celular busca el video que escribiste en RAM}; (3) \emph{aparece el video en la pantalla}; (4) \emph{el altavoz reproduce el sonido}; (5) \emph{escribes con el teclado un mensaje a tu mamá}; (6) \emph{el computador guarda el mensaje en el disco}; (7) \emph{la cámara del celular toma una foto}; (8) \emph{el procesador piensa cuánto demora el sol en salir mañana}; (9) \emph{la impresora saca una hoja con tu tarea}; (10) \emph{el computador recuerda tu contraseña guardada}. \emph{Tip:} algunas pueden parecer dos cosas; escoge la principal.
\end{softbox}

{\sffamily\bfseries\fontsize{8}{10}\selectfont\textcolor{milcNegro!55}{MARCA CUANDO LO HAYAS HECHO}}
\milccheckrow{Dibujé la tabla en el cuaderno.}
\milccheckrow{Clasifiqué las 10 acciones en una de las 4 funciones.}
\milccheckrow{No miré la siguiente actividad ni le copié al compañero.}

\begin{infoband}{milcVerde}


\textbf{Lo que va al cuaderno (Actividad 1 · IDENTIFICA).} Título: «Actividad 1 --- 10 acciones y sus funciones». \textbf{Formato:} tabla de 10 filas, 2 columnas. \textbf{Extensión:} 10 filas. \textbf{Sabes que terminaste cuando} cada una de las 10 acciones tiene su función asignada y puedes justificarla en voz alta.
\end{infoband}

\puente{Después del cuestionario, aprendes el nombre de cada función con la pieza del computador que la cumple.}

\milcestacion{milcTurquesa}{Estación 2 de 3 · Entender}

\begin{softbox}{milcTurquesa}{milcGris}{Lo que vas a entender}
Lo que acabaste de clasificar fueron \textbf{las 4 funciones} de cualquier máquina inteligente. Ahora vas a aprender el nombre técnico de cada función y \textbf{qué pieza del computador la cumple}. Esto es como el relojero: cada engranaje tiene su pieza, cada pieza tiene su función.
\end{softbox}

\milcpaso{1}{milcTurquesa}{\textbf{Función 1 --- ENTRADA (\emph{input} en inglés).} Es por donde llega la información al computador. \textbf{Piezas:} teclado, ratón (mouse), pantalla táctil, micrófono, cámara, escáner, joystick de juegos. \textbf{Cómo reconocerla:} si tú estás \emph{dándole algo} al computador (tocando, escribiendo, hablando, mostrándole), es entrada. \emph{Ejemplo:} cuando tocas la pantalla para abrir TikTok, esa es la entrada.}
\milcpaso{2}{milcTurquesa}{\textbf{Función 2 --- PROCESO.} Es donde el computador \emph{piensa}: hace cálculos, busca, decide, compara. \textbf{Piezas:} la CPU (\emph{unidad central de proceso}, es el ``cerebro'' del computador) y la GPU (la que se encarga de los gráficos pesados de los juegos). \textbf{Cómo reconocerla:} si la información se está \emph{transformando} en algo nuevo, eso es proceso. \emph{Ejemplo:} TikTok busca un video que coincide con lo que escribiste; ese trabajo de búsqueda es proceso.}
\milcpaso{3}{milcTurquesa}{\textbf{Función 3 --- ALMACENAMIENTO.} Es donde el computador \emph{guarda} la información. Hay 2 tipos. \textbf{RAM (memoria temporal):} guarda lo que estás usando AHORA. Si apagas el equipo, se borra. \textbf{Disco duro (memoria permanente):} guarda los archivos, las fotos, los videos, los programas. Aunque apagues el equipo, queda. \textbf{Cómo reconocerla:} si el computador \emph{recuerda} algo, eso es almacenamiento. \emph{Ejemplo:} la lista de tus videos vistos queda guardada (disco). Tu contraseña queda en RAM mientras la digitas.}
\milcpaso{4}{milcTurquesa}{\textbf{Función 4 --- SALIDA (\emph{output} en inglés).} Es como el computador \emph{te muestra} el resultado. \textbf{Piezas:} pantalla, altavoces, audífonos, impresora, vibrador del celular. \textbf{Cómo reconocerla:} si el computador \emph{te está mostrando algo}, eso es salida. \emph{Ejemplo:} cuando el video de TikTok aparece en pantalla y suena en el altavoz, eso es salida (visual + auditiva al mismo tiempo).}

\begin{drawbox}{milcTurquesa}{Las 4 funciones del computador}
\textbf{1. ENTRADA} (input) --- por donde llega la información. \\[1mm]
\textbf{2. PROCESO} --- donde el computador piensa (CPU). \\[1mm]
\textbf{3. ALMACENAMIENTO} --- donde guarda (RAM y disco duro). \\[1mm]
\textbf{4. SALIDA} (output) --- por donde muestra el resultado.
\end{drawbox}

\begin{softbox}{milcMostaza}{milcGris}{Errores comunes y su corrección}
\textbf{Confusiones frecuentes entre las 4 funciones.} (1) \emph{``La pantalla es solo salida''} --- Falso si es pantalla táctil: ahí es entrada (cuando tocas) Y salida (cuando muestra). (2) \emph{``El disco duro y la RAM son lo mismo''} --- Falso. RAM es corto plazo (se borra al apagar); disco duro es largo plazo (queda). (3) \emph{``El cerebro es el disco''} --- Falso. El cerebro es la CPU. El disco es solo memoria. (4) \emph{``El celular y el computador son distintos''} --- Falso. Tienen las mismas 4 funciones, solo varía el tamaño.
\end{softbox}

\begin{infoband}{milcTurquesa}


\textbf{Lo que va al cuaderno (Actividad 2 · EXPLICA).} Título: «Actividad 2 --- Las 4 funciones con sus piezas». \textbf{Formato:} tabla de 4 filas, 3 columnas. Columna 1: nombre de la función. Columna 2: 2 piezas que la cumplen. Columna 3: 1 ejemplo tuyo de tu día. \textbf{Extensión:} 4 filas. \textbf{Sabes que terminaste cuando} podrías mostrarle la tabla a un primo más pequeño y él entendería las 4 funciones en 2 minutos.
\end{infoband}

\puente{Con las 4 funciones claras, dibujas tu diagrama y trazas el recorrido completo de una acción real (jugar Roblox o escribir un correo).}

\milcestacion{milcMagenta}{Estación 3 de 3 · Hacer}

\begin{softbox}{milcMagenta}{milcGris}{Cómo vas a construir tu producto}
\textbf{Actividad 3 · APLICA --- Tu diagrama del computador en una hoja completa} (28 min · individual). Vas a dibujar en una hoja del cuaderno un \textbf{diagrama del computador} con las 4 funciones, las piezas que las cumplen, las flechas del flujo, y \textbf{un caso completo aplicado}: el recorrido de la información cuando juegas Roblox (o cuando ves un video, escoge tú).
\end{softbox}

\milcpaso{1}{milcMagenta}{\textbf{Paso 1 --- Título y formato.} En la parte superior de la hoja escribe: \textbf{``Mi diagrama del computador''} y debajo: \emph{``[Tu nombre], grado 6[tu letra], año 2026''}. Subraya el título. Pon la hoja en formato \emph{horizontal} (apaisada) porque vas a necesitar espacio a lo ancho.}
\milcpaso{2}{milcMagenta}{\textbf{Paso 2 --- Dibuja un rectángulo grande en el centro.} Adentro escribe \textbf{``COMPUTADOR''}. Divide el rectángulo en 4 cuadrantes (líneas horizontales y verticales en el medio). Etiqueta cada cuadrante con una función: arriba-izquierda \textbf{ENTRADA}, arriba-derecha \textbf{PROCESO}, abajo-derecha \textbf{ALMACENAMIENTO}, abajo-izquierda \textbf{SALIDA}.}
\milcpaso{3}{milcMagenta}{\textbf{Paso 3 --- Pon 2 piezas por función.} En cada cuadrante dibuja un ícono pequeño de \textbf{2 piezas que cumplen esa función}. \emph{Entrada:} teclado y pantalla táctil (o ratón y micrófono). \emph{Proceso:} CPU (un chip cuadrado) y GPU. \emph{Almacenamiento:} RAM (rectángulo delgado) y disco duro (cilindro). \emph{Salida:} pantalla y altavoz. Etiqueta cada ícono.}
\milcpaso{4}{milcMagenta}{\textbf{Paso 4 --- Flechas del flujo.} Con flechas curvas conecta: \textbf{Entrada → Proceso → Almacenamiento → Proceso → Salida}. Sí, el proceso aparece 2 veces porque la CPU trabaja antes y después de guardar. Si no entendieron del todo el flujo, así se aclara: la información entra, se procesa, se guarda lo que se necesita, se vuelve a procesar para mostrarla, y sale.}
\milcpaso{5}{milcMagenta}{\textbf{Paso 5 --- Caso aplicado: ``Juego Roblox''.} A la derecha del rectángulo grande, escribe un caso completo de 5 líneas. Una línea por función. \emph{Ejemplo:} \emph{(1) Toco la pantalla del celular para entrar a Roblox --- ENTRADA. (2) La CPU del celular busca el juego en el disco --- PROCESO. (3) El juego se carga en la RAM para usarse rápido --- ALMACENAMIENTO. (4) La CPU calcula los movimientos del personaje --- PROCESO. (5) La pantalla muestra el juego y el altavoz reproduce los sonidos --- SALIDA.}.}
\milcpaso{6}{milcMagenta}{\textbf{Paso 6 --- Firma de cierre.} En la esquina inferior derecha escribe: \emph{``Ahora entiendo qué pasa dentro del computador cuando lo uso.''} y firma con tu nombre completo y fecha. La idea es que esta hoja se vuelva un mapa de consulta cuando necesites entender un programa o un juego en el futuro.}

\begin{drawbox}{milcMagenta}{Antes de mostrar tu diagrama}
\checkbox El diagrama ocupa una hoja completa horizontal. \\[1mm]
\checkbox Están las 4 funciones etiquetadas en cuadrantes. \\[1mm]
\checkbox Cada función tiene 2 piezas dibujadas con su nombre. \\[1mm]
\checkbox Hay flechas que conectan las funciones (entrada → proceso → almacenamiento → proceso → salida). \\[1mm]
\checkbox Está el caso aplicado de 5 líneas (Roblox u otro). \\[1mm]
\checkbox Está firmado con nombre y fecha.
\end{drawbox}

\begin{softbox}{milcMagenta}{milcGris}{Plantilla de trabajo}
\textbf{Estructura del diagrama.} \textit{Centro:} rectángulo grande dividido en 4 cuadrantes con etiquetas. \textit{En cada cuadrante:} 2 piezas dibujadas. \textit{Flechas:} entrada → proceso → almacenamiento → proceso → salida. \textit{Lado derecho:} caso aplicado de 5 líneas. \textit{Esquina inferior:} frase de cierre + firma.
\end{softbox}

\begin{infoband}{milcMagenta}


\textbf{Lo que va al cuaderno (Actividad 3 · APLICA).} Título: «Actividad 3 --- Mi diagrama del computador». \textbf{Formato:} hoja completa horizontal con rectángulo de 4 cuadrantes, piezas dibujadas, flechas y caso aplicado. \textbf{Extensión:} 1 hoja entera. \textbf{Sabes que terminaste cuando} usando tu diagrama puedes explicar lo que pasa adentro del computador cuando juegas, escribes o ves un video.
\end{infoband}

\puente{El producto es el diagrama firmado con las 4 funciones, las piezas y el caso completo.}

\milcestacion{milcMostaza}{Producto de la sesión}
\textbf{Diagrama del computador con las 4 funciones · caso aplicado · firmado}.
\begin{softbox}{milcMostaza}{milcGris}{Criterios de calidad}
(1) El diagrama cubre la hoja completa horizontal. (2) Las 4 funciones están etiquetadas con sus piezas dibujadas. (3) Las flechas conectan las funciones en el orden correcto. (4) El caso aplicado tiene 5 líneas, una por función. (5) El diagrama está firmado y fechado.
\end{softbox}

\puente{Antes de salir, verificas que tu diagrama no tenga flechas perdidas: cada flecha conecta una función con la siguiente.}

\milcestacion{milcVino}{Evaluación liberadora · cinco dimensiones}

\begin{softbox}{milcVino}{milcGris}{Sentido de la evaluación}
Evaluar no es solo revisar una nota. Es reconocer qué aprendí, qué cambió en mí, qué emociones manejé, cómo me relaciono con la ciudad y la comunidad, y qué le dejaré a quien viene detrás.
\end{softbox}

\textbf{Mi reflexión liberadora en cinco dimensiones.} Completa cada frase iniciada:

\small
\begin{tabularx}{\textwidth}{>{\bfseries\RaggedRight\arraybackslash}p{3.4cm}Y>{\RaggedRight\arraybackslash}p{4.6cm}}
\rowcolor{milcVino}\color{white}Dimensión & \color{white}\bfseries Mi respuesta (frame starter) & \color{white}\bfseries Algo que descubrí\\
1. Personal & Lo que más cambió en mí fue \linead{6.5cm} & \linead{4.2cm}\\
2. Emocional & La emoción más fuerte fue \linead{2.5cm} y la regulé \linead{3cm} & \linead{4.2cm}\\
3. Ciudadana & En democracia esto importa porque \linead{6cm} & \linead{4.2cm}\\
4. Local (Cartago) & En mi barrio sería útil para \linead{6.5cm} & \linead{4.2cm}\\
5. Intergeneracional & A mi abuela/abuelo le diría \linead{6.5cm} & \linead{4.2cm}\\
\end{tabularx}
\normalsize

\begin{infoband}{milcVino}
\textbf{Para la próxima sustentación.} Vuelve a esta tabla en seis meses. Verás que las respuestas cambian — no porque lo aprendido cambie, sino porque tú cambias. La evaluación liberadora es una conversación contigo a través del tiempo.
\end{infoband}


\puente{Cerramos con 3 voces que te ayudan a entender por qué saber cómo funcionan las máquinas te hace más libre frente a ellas.}

\milcestacion{milcGradeDark}{Triángulo de pensamiento · cierre filosófico}

\begin{softbox}{milcVino}{milcGris}{Dussel · lente del nosotros}
\textit{Quien conoce su herramienta por dentro, la usa libre. Quien no la conoce, depende de quien sí.}

{\footnotesize\itshape\color{milcNegro!70}--- formulación de esta guía, al modo de Enrique Dussel. No es una cita textual.}

Hay 2 maneras de usar un computador: como \emph{mago} (\textbf{aprietas botones sin entender}, esperando que funcione) o como \emph{relojero} (\textbf{sabes qué hace cada parte}, entonces puedes resolver, ajustar, escoger). \textbf{La segunda manera te hace libre frente a la máquina.} La primera te hace dependiente. Hoy diste el primer paso para ser relojero: aprendiste las 4 funciones.

\textbf{¿Quiero seguir siendo mago con las máquinas (apretar y esperar), o quiero ser relojero (entender y decidir)?}
\end{softbox}
\linead{\linewidth}\\[1mm]
\linead{\linewidth}

\begin{softbox}{milcMostaza}{milcGris}{Estoicismo · Marco Aurelio · cuidado interior}
\textit{Pregunta siempre qué es la cosa en sí misma, no qué te parece que es. Esa pregunta separa al sabio del distraído.}

{\footnotesize\itshape\color{milcNegro!70}--- formulación de esta guía, al modo de Marco Aurelio. No es una cita textual.}

Cuando ves un computador, lo fácil es pensar \emph{``es un aparato que sirve para esto''}. Lo sabio es preguntar: \emph{``¿qué es por dentro?''}. Lo mismo con un celular, una nevera inteligente, un cajero electrónico. Detrás de cualquier máquina hay esas mismas 4 funciones. \textbf{Saber esto te entrena a no asustarte con tecnología nueva}: si tiene esas 4 funciones, las puedes entender.

\textbf{¿Me conformo con saber qué \emph{hace} una cosa, o también me pregunto qué \emph{es} por dentro?}
\end{softbox}
\linead{\linewidth}\\[1mm]
\linead{\linewidth}

\begin{softbox}{milcTurquesa}{milcGris}{Floridi · lente de la infoesfera}
\textit{Vivimos en una sociedad de máquinas inteligentes. Quien no entiende cómo funcionan, las usa sin ser su dueño.}

{\footnotesize\itshape\color{milcNegro!70}--- formulación de esta guía, al modo de Luciano Floridi. No es una cita textual.}

En 20 años habrá máquinas que ni se han inventado. Pero \textbf{las 4 funciones siguen siendo las mismas}: entrada, proceso, almacenamiento, salida. Si dominas este modelo, dominas cualquier máquina futura. Esto no es teoría: es preparación para una vida entera de tecnologías nuevas.

\textbf{¿Estoy aprendiendo lo de hoy o las ideas que sirven para 30 años?}
\end{softbox}
\linead{\linewidth}\\[1mm]
\linead{\linewidth}

\begin{infoband}{milcNegro}
\textbf{Nota para el docente.} Las tres formulaciones son del autor de la guía, no citas textuales de los pensadores: recogen su lente para pensar el tema, no sus palabras. Si vas a citarlos en clase, remite a la obra original.
\end{infoband}

\milcestacion{milcVino}{Mi autoevaluación · marca una casilla por fila}

{\small
\setlength{\extrarowheight}{2.2pt}
\begin{tabularx}{\textwidth}{>{\RaggedRight\arraybackslash\bfseries}p{3.0cm}YYY}
\rowcolor{milcVino}\color{white}\bfseries Criterio & \color{white}\bfseries Lo logré & \color{white}\bfseries En proceso & \color{white}\bfseries Necesito apoyo\\[.6mm]
Comprensión del tema & \milcopcion{Explico las ideas con mis palabras.} & \milcopcion{Reconozco las ideas, falta precisión.} & \milcopcion{Necesito repasar lo básico.}\\[.8mm]
\rowcolor{milcGris}Pensamiento computacional & \milcopcion{Usé los pilares al armar mi producto.} & \milcopcion{Usé algunos pilares.} & \milcopcion{Necesito releer la estación 2.}\\[.8mm]
Aplicación y producto & \milcopcion{Mi producto cumple los criterios.} & \milcopcion{Está armado, con ajustes pendientes.} & \milcopcion{Aún está en construcción.}\\[.8mm]
\rowcolor{milcGris}Reflexión 5 dimensiones & \milcopcion{Respondí cada una con honestidad.} & \milcopcion{Respondí algunas con detalle.} & \milcopcion{Mis respuestas son muy generales.}\\[.8mm]
Triángulo de pensamiento & \milcopcion{Conecté las 3 voces con mi trabajo.} & \milcopcion{Conecté al menos una.} & \milcopcion{No las relaciono con mi proceso.}\\[.8mm]
\end{tabularx}}

\vspace{3mm}
\textbf{Hoy entendí que:}\\[1mm]
\linead{\linewidth}\\[1mm]
\linead{\linewidth}

\textbf{Mi compromiso para la próxima sesión es:}\\[1mm]
\linead{\linewidth}\\[1mm]
\linead{\linewidth}

\begin{softbox}{milcMostaza}{milcGris}{Cita de despedida · Marco Aurelio}
\textit{``El obstáculo en el camino se vuelve el camino. Lo que estorba al camino se vuelve camino.''}\\[1mm]
Cada error de hoy, bien observado, se convierte en pieza del camino que pisarás mañana.
\end{softbox}

\small
\begin{tabularx}{\textwidth}{>{\bfseries}p{3.6cm}Y}
\rowcolor{milcGradeDark!12}Nombre completo: & \linead{10cm}\\
Fecha: & \linead{4cm} \quad Curso/grupo: \linead{4cm}\\
Mi firma: & \linead{9cm}\\
\end{tabularx}
\normalsize

\begin{infoband}{milcGradeDark}
\textbf{Para la próxima sesión.} Esta semana, cuando uses tu celular, juega a identificar las 4 funciones. Por ejemplo, al ver un TikTok: ¿qué es entrada? ¿qué es proceso? ¿qué es salida? Si lo haces 3 veces, ya lo tendrás claro. La próxima clase abrimos el gabinete del computador y vemos las piezas internas: CPU, RAM y disco duro por dentro.
\end{infoband}

\end{document}
